%% ma: L11-B16-0003
%% lop: 11
%% chuong: 5
%% bai: 16
%% dang: 11.16.2
%% loai: TL
%% muc: VD
%% dap_an: "-"
%% nguon: "tai-lieu-hoc-tap-mon-toan-lop-11-theo-chuong-trinh-giao-duc-pho-thong-2018.pdf (Lớp Toán Thầy Duy, Phạm Lê Duy, THPT Chu Văn An), Chương 3 Bài 2, Dạng 1, Ví dụ 3 (tr.351)"
%% kien_thuc: "Lớp 11 Bài 16 (giới hạn của hàm số)"
%% trang_thai: cho-duyet
%% ly_do: "Dạng mới đề xuất (11.16.x), chờ giáo viên duyệt dạng; tài liệu gốc không có lời giải mẫu, đáp án do kho tự giải và kiểm bằng sympy."
%% ghi_chu: "Ý f) của đề gốc, lim_{x->1} (sqrt(2x+7)+x+4)/(x^3-4x^2+3), có tử bằng 8 và mẫu bằng 0 nên không có giới hạn (hai phía trái dấu); nghi đề gõ sai dấu (nếu là x-4 thì giới hạn = -4/15), nên kho bỏ ý f)."
\begin{ex}
Tìm giới hạn

a) $\lim\limits_{x\to0}\dfrac{\sqrt{1+2x}-1}{2x}$;

b) $\lim\limits_{x\to0}\dfrac{4x}{\sqrt{9+x}-3}$;

c) $\lim\limits_{x\to1}\dfrac{\sqrt{2x+7}-3}{2-\sqrt{x+3}}$;

d) $\lim\limits_{x\to0}\dfrac{\sqrt{1+x^2}-1}{\sqrt{x^2+16}-4}$;

e) $\lim\limits_{x\to2}\dfrac{\sqrt{x+2}-\sqrt{2x}}{\sqrt{x-1}-\sqrt{3-x}}$.
\loigiai{\sloppy a) \[\dfrac{\sqrt{1+2x}-1}{2x}=\dfrac{2x}{2x\left(\sqrt{1+2x}+1\right)}=\dfrac1{\sqrt{1+2x}+1}\to\dfrac12.\]

b) \[\dfrac{4x}{\sqrt{9+x}-3}=\dfrac{4x\left(\sqrt{9+x}+3\right)}{x}=4\left(\sqrt{9+x}+3\right)\to24.\]

c) Nhân liên hợp tử và mẫu: \[\begin{aligned}\dfrac{\sqrt{2x+7}-3}{2-\sqrt{x+3}}=\dfrac{2(x-1)\left(2+\sqrt{x+3}\right)}{-(x-1)\left(\sqrt{2x+7}+3\right)}\to-\dfrac{2\cdot4}{6}\\ &=-\dfrac43.\end{aligned}\]

d) \[\dfrac{\sqrt{1+x^2}-1}{\sqrt{x^2+16}-4}=\dfrac{x^2\left(\sqrt{x^2+16}+4\right)}{x^2\left(\sqrt{1+x^2}+1\right)}\to\dfrac82=4.\]

e) Nhân liên hợp: \[\begin{aligned}\dfrac{(2-x)\left(\sqrt{x-1}+\sqrt{3-x}\right)}{(2x-4)\left(\sqrt{x+2}+\sqrt{2x}\right)}=\dfrac{-\left(\sqrt{x-1}+\sqrt{3-x}\right)}{2\left(\sqrt{x+2}+\sqrt{2x}\right)}\to-\dfrac{2}{2\cdot4}\\ &=-\dfrac14.\end{aligned}\]}
\end{ex}
